% !TeX program = xelatex
% !TEX root = main.tex
% !TeX encoding = UTF-8
\documentclass[10.5pt,compsoc]{CjC}
%\usepackage{CJKutf8}
%\usepackage{CJK}
\usepackage{graphicx}
\usepackage{footmisc}
\usepackage{subfigure}
\usepackage{url}
\usepackage{multirow}
\usepackage[noadjust]{cite}
\usepackage{amsmath,amsthm}
\usepackage{amssymb,amsfonts}
\usepackage{booktabs}
\usepackage{color}
\usepackage{ccaption}
\usepackage{float}
\usepackage{fancyhdr}
\usepackage{caption}
\usepackage{xcolor,stfloats}
\usepackage{comment}
\setcounter{page}{1}
\graphicspath{{figures/}}
\usepackage{cuted}
\usepackage{captionhack}
\usepackage{epstopdf}

%===============================%
% XeLaTeX + CJK font setup. CjC.cls is based on ctexart; these explicit fonts
% reduce the risk of garbled Chinese output on Windows TeX Live.
\setmainfont{Times New Roman}
\setCJKmainfont{SimSun}
\setCJKsansfont{SimHei}
\setCJKmonofont{FangSong}

% Manual body and heading spacing overrides for the compact CjC two-column
% layout. These settings intentionally override the class defaults so that long
% Chinese section headings never visually collide with adjacent paragraphs.
\setlength{\parskip}{0.38\baselineskip plus 0.08\baselineskip minus 0.04\baselineskip}
\makeatletter
\renewcommand\section{\@startsection {section}{1}{\z@}%
  {4.0ex \@plus 1.0ex \@minus .2ex}%
  {2.4ex \@plus .5ex}%
  {\normalfont\zihao{4-}\heiti\bfseries\raggedright}}
\renewcommand\subsection{\@startsection{subsection}{2}{\z@}%
  {3.2ex \@plus .8ex \@minus .2ex}%
  {1.8ex \@plus .4ex}%
  {\normalfont\zihao{5}\heiti\bfseries\raggedright}}
\renewcommand\subsubsection{\@startsection{subsubsection}{3}{\z@}%
  {2.6ex \@plus .6ex \@minus .2ex}%
  {1.4ex \@plus .3ex}%
  {\normalfont\zihao{5}\songti\bfseries\raggedright}}
\makeatother

\headevenname{\mbox{\quad} \hfill  \mbox{\zihao{-5}{\songti  计\quad \quad 算\quad \quad 机\quad \quad 学\quad \quad 报 } \hspace {50mm} \mbox{\songti  20?? 年 }}}%
\headoddname{\songti  ? 期 \hfill
课程论文作者：多模态大模型在智能电网中的应用综述 }%

\renewcommand{\thefootnote}{\fnsymbol{footnote}}
\setcounter{footnote}{0}
\renewcommand\footnotelayout{\zihao{5-}}

\newtheoremstyle{mystyle}{0pt}{0pt}{\normalfont}{1em}{\bf}{}{1em}{}
\theoremstyle{mystyle}
\renewcommand\figurename{figure~}
\renewcommand{\thesubfigure}{(\alph{subfigure})}
\newcommand{\upcite}[1]{\textsuperscript{\cite{#1}}}
\renewcommand{\labelenumi}{(\arabic{enumi})}
\newcommand{\tabincell}[2]{\begin{tabular}{@{}#1@{}}#2\end{tabular}}
\newcommand{\abc}{\color{white}\vrule width 2pt}
\makeatletter
\renewcommand{\@biblabel}[1]{[#1]\hfill}
\makeatother
\setlength\parindent{2em}

\begin{document}
\hyphenpenalty=50000
\makeatletter
\newcommand\mysmall{\@setfontsize\mysmall{7}{9.5}}
\newenvironment{tablehere}
  {\def\@captype{table}}

\let\temp\footnote
\renewcommand \footnote[1]{\temp{\zihao{-5}#1}}

\thispagestyle{plain}%
\thispagestyle{empty}%
\pagestyle{CjCheadings}

\begin{table*}[!t]
\vspace {-13mm}
\begin{tabular}{p{168mm}}
\zihao{5-}
\songti
第??卷\quad 第?期 \hfill 计\quad 算\quad 机\quad 学\quad 报\hfill Vol. ??  No. ?
\zihao{5-}\songti
20??年?月 \hfill CHINESE JOURNAL OF COMPUTERS \hfill ???. 20?? \\
\hline\\[-4.5mm]
\hline\end{tabular}

\centering
\vspace{11mm}
\heiti
{\zihao{2} 多模态大模型在智能电网中的应用综述}

\vskip 5mm

{\zihao{3} \fangsong
课程论文作者$^{1)}$
}

\vspace{5mm}
\zihao{6}{\songti
$^{1)}$(课程论文，单位信息略)
}

\zihao{6}{\heiti
本文为课程论文初稿，作者、单位、基金及联系方式信息略。
}

\vskip 5mm
{\centering
\begin{tabular}{p{160mm}}
\zihao{5-}{
\setlength{\baselineskip}{16pt}\selectfont{
\noindent\heiti 摘\quad 要\quad   \songti
新型电力系统和智能电网在源网荷储协同、设备状态感知、调度运行和安全治理中产生了图像、视频、红外、文本、时序、拓扑和知识图谱等多源异构数据，传统单模态智能方法难以支撑复杂业务场景下的统一理解、交互推理和可信决策。多模态大模型通过视觉语言对齐、跨模态表示学习、指令微调、检索增强生成和知识图谱等机制，为电力数据融合与业务知识利用提供了新的技术路径。本文围绕多模态大模型在智能电网中的应用，梳理关键技术基础，构建多模态数据与任务体系，总结电力巡检、故障诊断、设备运维、调度辅助、应急演练和安全监管等典型应用，分析知识增强可信推理、工程部署、数据治理、可靠性、实时性、可解释性和安全边界等挑战，并展望电力专用多模态基础模型、物理约束融合、动态知识增强、人在回路决策和多智能体协同等方向。
\par}}\\[2mm]

\zihao{5-}{\noindent
\heiti 关键词  \quad \songti  多模态大模型；智能电网；视觉语言模型；电力巡检；故障诊断；调度决策；知识增强
}\\[2mm]
\zihao{5-}{\heiti 中图法分类号 	\songti
TP391 \rm{\quad \quad \quad     }
\heiti DOI号: \songti
投稿时不提供DOI号}
\end{tabular}}

\vskip 7mm

\begin{center}
\zihao{3}{ {\heiti A Survey on Applications of Multimodal Large Models in Smart Grids}}\\
\vspace {5mm}
\zihao{5}{ {\heiti Course Paper Author$^{1)}$}}\\
\vspace {2mm}
\zihao{6}{\heiti {$^{1)}$(Course Paper, affiliation omitted)}}

\end{center}

\begin{tabular}{p{160mm}}
\zihao{5}{
\setlength{\baselineskip}{18pt}\selectfont{
{\bf Abstract}\quad
Smart grids generate heterogeneous data from inspection images, infrared sensing, operation logs, dispatching instructions, time series measurements, topology models, equipment ledgers, and domain knowledge bases. Conventional single modal artificial intelligence methods are often insufficient for unified perception, semantic understanding, interactive reasoning, and trustworthy decision support. This paper surveys multimodal large models for smart grids. It reviews the technical foundations, organizes smart grid data and tasks into a multimodal taxonomy, summarizes representative applications in inspection, defect recognition, fault diagnosis, operation and maintenance, dispatching assistance, emergency rehearsal, and safety supervision, and analyzes challenges in knowledge enhanced reasoning, engineering deployment, data governance, reliability, real time constraints, explainability, and security boundaries. Future directions include power domain multimodal foundation models, physical constraint integration, dynamic knowledge enhancement, human in the loop decision making, and multi agent collaboration.
\par}}\\

\setlength{\baselineskip}{18pt}\selectfont{
\zihao{5}{\noindent
\vspace {5mm}
{\bf Keywords}\quad multimodal large models; smart grid; vision-language models; power inspection; fault diagnosis; dispatching decision; knowledge enhancement
}\par}
\end{tabular}

\setlength{\tabcolsep}{2pt}
\begin{tabular}{p{0.05cm}p{16.15cm}}
\multicolumn{2}{l}{\rule[4mm]{40mm}{0.1mm}}\\[-3mm]
&\begin{songti}
课程论文初稿：作者、单位、基金、邮箱及手机号等投稿专用信息略。
\end{songti}
\end{tabular}\end{table*}
\clearpage\clearpage
\begin{strip}
\vspace {-13mm}
\end{strip}
\heiti
\zihao{5}
\setlength{\baselineskip}{17pt}\selectfont
\vskip 3mm


\section{引言}

\songti
新型电力系统以高比例新能源接入、源网荷储协同、分布式资源互动和数字化调度为重要特征。与传统电力系统相比，智能电网不仅要处理潮流、负荷、状态量等结构化数据，还要面对巡检图像、红外图像、作业视频、缺陷文本、设备台账、调度规程、故障工单和拓扑图结构等多源异构信息。已有研究表明，大语言模型正在进入电力系统规划、运行、市场、运维和网络安全等环节\cite{ghafari2025comprehensive,mirshekali2025review,niu2025llm_power}，但智能电网的真实业务对象往往并非纯文本任务，而是由多模态感知、领域知识检索、状态推理和业务流程协同共同构成。

多模态大模型(multimodal large model, MLLM)为上述问题提供了新的研究路径。通用大模型的发展以 Transformer 架构\cite{vaswani2017attention}和大语言模型的上下文学习能力\cite{brown2020language}为基础，视觉语言模型进一步通过图文对齐与视觉指令微调将图像理解、自然语言交互和推理解释连接起来\cite{radford2021learning,liu2023visual}。在电力领域，这类能力有望将巡检图像识别、缺陷描述、规程查询、故障处置和调度辅助从割裂的单点模型推进为统一的人机协同智能系统\cite{li2024foundation_power,xiang2025mllm_power}。

然而，多模态大模型不能被简单理解为通用模型在电力数据上的直接迁移。电力系统具有强安全、强实时、强物理约束和强责任边界等特点，模型输出需要满足设备约束、运行规程和可审计要求。现有综述已对电力系统大语言模型、能源系统大模型和电力大模型关键问题进行了初步梳理\cite{sarwar2025large,fan2025cognitive,zhang2026large,du2026key}，但对``多模态数据类型--模型机制--电网任务--可信约束''之间关系的系统分析仍显不足。为此，本文围绕多模态大模型在智能电网中的应用展开综述，重点回答三个问题：智能电网中哪些数据和任务需要多模态大模型；多模态大模型在巡检、诊断、调度和知识增强中如何发挥作用；面向电力生产系统部署时还存在哪些可信、安全和工程挑战。

本文余下部分组织如下：第2节介绍多模态大模型关键技术基础；第3节构建智能电网多模态数据与任务体系；第4节总结典型应用场景；第5节讨论知识增强、RAG与可信推理；第6节分析工程部署、安全可信和评测挑战；第7节展望未来研究方向；第8节给出结论。


\section{多模态大模型关键技术基础}


\subsection{Transformer与大语言模型}

Transformer 通过自注意力机制实现序列建模能力，成为当前大语言模型和多模态大模型的核心架构基础\cite{vaswani2017attention}。大语言模型在规模化预训练后表现出上下文学习、少样本迁移、指令跟随和自然语言生成能力\cite{brown2020language}，使其能够处理电力规程问答、运维报告生成、故障处置建议和调度语义理解等知识密集型任务\cite{ghafari2025comprehensive,niu2025llm_power}。但通用语言模型的训练语料并不天然覆盖电力设备、调度规程和运行约束，因此在电力系统中使用时需要领域知识注入、数据安全隔离和输出校验。


\subsection{视觉语言模型与多模态指令微调}

视觉语言模型通过图像和文本之间的语义对齐，使模型具备跨模态表示和零样本视觉理解能力。CLIP 通过自然语言监督学习可迁移视觉表示\cite{radford2021learning}，LLaVA 等视觉指令微调方法进一步使模型具备图像问答、描述生成和交互式解释能力\cite{liu2023visual}。在电力巡检中，图像、红外、设备名称、缺陷类型和巡检规程之间存在天然的语义对应关系，因此视觉语言模型可用于缺陷识别、异常解释和巡检报告生成。PowerVLM 等电力视觉语言模型研究已开始探索联邦学习、模型剪枝和行业场景适配\cite{ouyang2026powervlm}。


\subsection{检索增强生成与知识图谱}

检索增强生成(retrieval augmented generation, RAG)通过将生成模型与外部文档检索结合，使模型输出可以依托非参数化知识来源\cite{lewis2020retrieval}。电力系统中的规程制度、设备台账、历史工单和应急预案具有明确版本与责任属性，RAG 可用于提高回答可追溯性并降低幻觉风险。知识图谱则能够以结构化方式表达设备、缺陷、部件、工况、规程和处置动作之间的关系\cite{ji2022survey}，为多模态大模型提供可查询、可推理的领域知识空间。近期 HG-RAG 等框架也表明，层次化图结构可进一步增强复杂知识检索和生成能力\cite{shen2026hgrag}。


\subsection{电力领域适配技术}

电力领域适配涉及模型能力、数据约束和部署环境三方面。首先，电力专用指令微调和任务数据集可提高模型对设备缺陷、故障现象和调度术语的理解能力。其次，电力数据存在隐私和安全限制，联邦学习、私有化部署、数据脱敏和权限控制是行业模型落地的重要条件\cite{ouyang2026powervlm,du2026key}。再次，智能电网边缘侧设备和现场网络环境有限，模型压缩、边云协同和电力智算基础设施将影响多模态大模型的工程可用性\cite{fangfang2025computing}。


\section{智能电网多模态数据与任务体系}


\subsection{多模态数据类型}

智能电网的多模态数据可分为七类：一是可见光图像和视频，用于设备外观、线路通道和作业现场识别；二是红外和局部放电等状态感知数据，用于反映设备热异常或绝缘缺陷；三是电压、电流、功率、频率等时序量测数据；四是电网拓扑、设备连接和站线关系等图结构数据；五是规程、标准、调度指令、检修票和工单等文本数据；六是设备台账、缺陷库和检修记录等业务数据；七是由上述数据抽象形成的电力知识图谱。不同模态在采集频率、噪声水平、语义粒度和安全等级上存在差异，决定了多模态融合不能仅靠简单拼接，而需要结合任务目标和业务约束进行建模\cite{li2024foundation_power,du2026key,zhengxiang2025data}。


\subsection{多模态任务谱系}

从任务层次看，智能电网多模态大模型可支持感知识别、语义理解、知识检索、状态推理、方案生成和人机协同六类能力。感知识别主要面向巡检图像缺陷检测和作业现场识别；语义理解关注图像、文本和设备台账之间的关系解释；知识检索用于规程问答和证据定位；状态推理面向故障诊断、运行态势研判和异常原因分析；方案生成用于抢修方案、巡检报告和应急演练脚本；人机协同则强调模型建议经人工复核后进入业务闭环\cite{ghafari2025comprehensive,mirshekali2025review,xiang2025mllm_power}。


\subsection{面向智能电网的分类框架}

本文建议从``数据模态--模型机制--电力任务--可信约束''四个维度理解多模态大模型应用。数据模态刻画输入来源，模型机制包括视觉语言对齐、指令微调、RAG、知识图谱和工具调用，电力任务对应巡检、诊断、调度、应急和运维等业务场景，可信约束则覆盖物理可行性、安全边界、实时性、可解释性和审计追踪。该框架有助于避免将多模态大模型泛化为单纯的图像识别或文本问答技术，也有助于区分不同场景的成熟度与风险等级。

\begin{figure*}[htbp]
\centering
\fbox{\begin{minipage}{0.92\textwidth}
\centering
\begin{tabular}{c}
\fbox{数据层：图像、视频、红外、文本、时序、拓扑、知识图谱}\\[1mm]
$\Downarrow$\\[-1mm]
\fbox{模型层：视觉语言对齐、指令微调、跨模态表示、工具调用}\\[1mm]
$\Downarrow$\\[-1mm]
\fbox{知识层：规程库、设备台账、历史工单、RAG、知识图谱}\\[1mm]
$\Downarrow$\\[-1mm]
\fbox{应用层：巡检识别、故障诊断、调度辅助、应急演练}\\[1mm]
$\Downarrow$\\[-1mm]
\fbox{可信治理层：物理约束、安全校验、证据追溯、人在回路}
\end{tabular}
\end{minipage}}
\caption{多模态大模型赋能智能电网的总体框架}
\label{fig:framework}
\end{figure*}


\subsection{与传统电力人工智能方法的关系}

多模态大模型并不替代传统机器学习、优化算法和电力机理模型，而是对其形成补充。传统方法在潮流计算、状态估计、故障分类和优化调度中具有明确数学结构和可验证性；多模态大模型则更适合开放语义理解、复杂文档交互、多源证据融合和业务流程协同。对于高安全等级的调度控制任务，应将模型输出定位为辅助建议，并通过物理模型、仿真平台和调度规则进行约束校验\cite{zhangjun2023operation,li2024risks}。


\section{多模态大模型在智能电网中的典型应用}


\subsection{电力巡检与设备缺陷识别}

电力巡检是多模态大模型较容易形成可见业务价值的应用场景，但其可用性并不等同于端到端自动判缺。传统视觉模型主要面向设备目标检测、缺陷分类和异常区域定位，多模态大模型则可进一步将可见光图像、红外图像、巡检视频、设备台账和运维规程组织为统一的语义解释过程。例如，电力视觉语言模型可在识别绝缘子破损、导线异物、表计异常和设备外观缺陷的基础上，生成包含缺陷位置、缺陷类别、相关规程和处置建议的初步报告\cite{ouyang2026powervlm,wang2025inspection,liu2026defect}。因此，巡检场景中 MLLM 的合理定位不是替代现场判定，而是在图像识别、证据组织和报告草拟环节提供辅助。

与表1中的``电力巡检''场景相对应，该类应用的主要输入模态包括图像、视频、红外、设备台账和巡检规程，典型机制包括视觉语言模型和指令微调。其优势在于能够把缺陷识别与业务语言解释连接起来，使巡检结果更容易进入工单系统和缺陷闭环治理。然而，该优势受到样本不均衡、跨场站泛化和复杂环境鲁棒性的制约。不同设备厂家、拍摄角度、光照天气和缺陷等级会显著改变图像分布，红外证据与可见光证据之间也可能存在时空对齐误差。若模型在少样本缺陷上过度归纳，或在模糊图像中给出过度确定的缺陷描述，便可能造成误报、漏报和不恰当的检修优先级排序。因此，后续研究应将巡检 MLLM 评价从单一识别准确率扩展到跨场景鲁棒性、证据一致性、缺陷等级校准、人工复核效率和工单闭环质量等指标。


\subsection{故障诊断、设备运维与抢修方案生成}

故障诊断和设备运维要求模型从多源证据中推断故障类型、故障位置和处置策略。与单纯分类任务不同，电力故障处置通常包含现场现象描述、保护动作记录、设备状态量、历史工单、图像证据和规程条款。大模型指令微调可用于生成抢修方案和运维建议\cite{deng2026repair}，类 DeepSeek 范式等推理模型也被用于探索电力装备故障诊断中的推理链组织\cite{jiang2025deepseek}。已有故障诊断综述指出，机器学习和大语言模型为新型电力系统故障诊断提供了新的方法补充，但仍需与机理知识和专家经验结合\cite{zhangjianliang2025fault}。

从工程角度看，故障诊断类应用尤其需要证据链。模型不应只给出``可能是某故障''的结论，还应明确依据哪些量测、图像、规程或历史案例，并给出不确定性提示。对于抢修方案生成，应将模型输出限制为辅助建议，经现场人员和调度人员审核后执行，避免生成式模型直接触发高风险操作。


\subsection{调度控制、运行状态感知与辅助决策}

调度控制和运行状态感知体现了多模态大模型在智能电网中的高价值与高风险并存特征。现有研究开始评估大语言模型在电力系统调度任务中的能力\cite{mongaillard2024scheduling}，也有工作探索将大语言模型作为配电网电压控制中的操作员角色\cite{yang2025operator}。PowerGraph-LLM 等图增强方法表明，将电网拓扑结构与语言模型结合，可能缓解通用 LLM 对节点、支路和运行状态关系表达不足的问题\cite{bernier2025powergraph}。中文研究也从运行控制、推理大模型和配电系统专用大模型角度讨论了智能决策前景\cite{zhangjun2023operation,yutao2026reasoning}。

与表1中的``调度辅助''场景相对应，调度类应用的输入通常同时包含时序量测、拓扑结构、告警日志、调度指令和规程文本，典型机制包括图增强 LLM、工具调用和约束校验。MLLM 在该场景中的功能应被限定为四类辅助环节：一是对运行日志、告警信息和多源量测进行态势摘要；二是检索并解释历史预案、调度规程和运行约束；三是生成候选调度方案或处置思路；四是在潮流计算、优化调度和仿真验证之后，对校验结果进行自然语言解释。上述流程中，模型生成的方案只能作为待验证候选项，不能越过潮流约束、设备限值、保护定值和调度权限进入控制链路。

调度场景的主要风险并非仅来自语言模型的幻觉，还来自物理约束被遗漏、拓扑状态被误读、实时数据滞后和操作权限边界不清。若模型建议导致节点电压越限、线路潮流越限或功率平衡破坏，即使其语义解释看似合理，也不具备工程可执行性。因此，当前阶段更稳妥的研究方向是将 MLLM 嵌入调度辅助和方案解释环节，并通过仿真平台、规则系统和人工调度员形成多级校验，而不是将其表述为自主控制器。


\subsection{应急演练、安全监管与工程实践}

应急演练和安全监管属于知识密集、流程密集但相对可控的应用场景。基于高级 RAG 与多模态大模型的电力应急演练研究表明，模型可将应急预案、现场图像、处置流程和问答交互结合起来，生成演练脚本、风险提示和处置建议\cite{liwenchao2025rag}。在电力工程实践中，行业大模型已被用于智能巡检、安全督察、调度辅助和知识服务等场景\cite{ghafari2025comprehensive,li2024foundation_power}。

这类场景的优势在于可先以离线或半在线方式部署，降低直接控制风险。其主要挑战是知识库版本管理、业务流程一致性和输出责任归属。若模型引用过期规程或错误工单，可能造成误导。因此，应急与监管类应用也需要建立证据追溯、版本控制和人工确认机制。

\begin{table*}[htbp]
\centering
\heiti 表1\quad 多模态大模型在智能电网典型应用中的任务与约束
\vspace{-2mm}
\begin{tabular}{p{22mm}p{31mm}p{32mm}p{36mm}p{33mm}}
\toprule
应用场景 & 主要输入模态 & 典型模型机制 & 主要优势 & 主要约束 \\
\midrule
电力巡检 & 图像、视频、红外、台账、规程 & 视觉语言模型、指令微调 & 缺陷候选识别、证据解释、报告草拟 & 样本不均衡、跨场景泛化、误报漏报 \\
故障诊断 & 状态量、工单、图像、历史案例 & RAG、推理链、知识图谱 & 多源证据汇聚、诊断依据组织、方案草案生成 & 幻觉风险、证据冲突、责任边界 \\
调度辅助 & 时序、拓扑、调度指令、规程 & 图增强LLM、工具调用、约束校验 & 态势摘要、预案检索、候选方案解释 & 实时性、物理约束、越权控制风险 \\
应急演练 & 预案、现场图像、流程文本 & RAG、多模态问答 & 情景推演、流程协同、风险提示 & 知识版本、流程一致性、人工复核 \\
\bottomrule
\end{tabular}
\label{tab:applications}
\end{table*}


\section{知识增强、RAG与电力领域可信推理}


\subsection{电力知识图谱与多模态知识组织}

电力知识图谱能够将设备、部件、缺陷、规程、状态量、检修记录和运行场景组织为结构化知识网络。与纯文本知识库相比，知识图谱能够显式表达实体关系和业务约束，有利于模型进行可解释推理\cite{ji2022survey}。在变电站智能运检、电力领域多模态知识图谱和能源数字网络知识图谱等研究中，知识图谱被用于支撑设备状态理解、跨模态检索和知识服务\cite{fuxuhui2025mmkg,zhouaihua2026kg}。

多模态知识图谱的关键难点在于实体对齐和动态更新。图像中的设备部件、台账中的设备编码、规程中的专业术语和工单中的故障描述可能对应同一实体或状态。如果对齐错误，后续检索和推理都会受到影响。因此，电力知识图谱建设需要结合标准化命名、设备编码体系、图像标注和业务流程规则。


\subsection{RAG在电力规程和业务知识中的应用}

RAG 将外部知识检索与模型生成结合，适合电力规程问答、故障处置建议和应急演练等知识密集型任务\cite{lewis2020retrieval,liwenchao2025rag}。在电力场景中，RAG 的评价重点不应只是答案是否流畅，还应关注检索证据是否权威、版本是否有效、引用是否可追溯以及不同证据之间是否冲突。层次化图检索等方法可进一步提高复杂知识结构下的检索组织能力\cite{shen2026hgrag}。


\subsection{可信推理机制}

可信推理要求模型输出从``生成文本''转向``给出可验证建议''。在故障诊断中，模型应说明使用了哪些图像证据、状态量、历史工单和规程条款；在调度辅助中，模型应说明建议是否经过潮流约束、设备限值和安全规则校验；在应急演练中，模型应给出处置流程与预案来源。图增强语言模型、知识图谱和 RAG 可以共同构成证据链，为多模态大模型提供可审计的推理基础\cite{bernier2025powergraph,ji2022survey,li2024risks}。


\section{工程部署、安全可信与评测挑战}


\subsection{数据治理与隐私安全}

电力数据具有关键基础设施属性，涉及企业生产、用户隐私和电网安全。多模态大模型训练和部署面临数据孤岛、标注成本高、数据质量不一致和共享限制等问题\cite{zhengxiang2025data,du2026key}。联邦学习、私有化部署、模型压缩和权限控制可在一定程度上缓解数据安全与模型能力之间的矛盾\cite{ouyang2026powervlm}，但也会带来模型更新、性能评估和系统运维的新问题。


\subsection{可靠性、实时性与物理约束}

智能电网中的模型输出必须同时满足语义合理性与物理可行性。大模型可能生成表述完整、逻辑连贯但工程上不可执行的建议，例如忽略网络拓扑变化、违反设备热稳定限值、造成节点电压越限，或将不适用于当前运行方式的调度规程错误迁移到实时场景。智能电网大模型风险研究指出，模型在电网环境中的误用可能导致安全、隐私和运行风险\cite{li2024risks}。因此，表2中将``物理约束''列为影响调度控制和故障处置的核心挑战，其本质是生成式建议与可验证电力计算之间的匹配问题。

面向调度控制和故障处置，MLLM 输出应首先转化为结构化候选动作，而不是直接作为自然语言建议进入业务系统。该候选动作至少应包含对象设备、操作类型、预期目标、适用工况、证据来源和不确定性提示。随后，系统应依次执行静态规则校验、潮流可行性校验、仿真验证和人工确认。静态规则校验用于排除违反调度规程、权限边界和设备限值的动作；潮流计算用于检查候选动作是否导致节点电压越限、线路潮流越限或功率平衡破坏；动态仿真或故障仿真用于评估扰动条件下的安全裕度；人工确认则记录调度员对模型建议、校验结果和最终操作的责任判断。只有通过上述环节的建议，才可作为辅助决策依据；未通过任一环节的建议应被回退为不可执行解释。

实时性约束进一步限制了 MLLM 在在线控制中的角色。对于秒级或毫秒级控制任务，模型生成、检索和校验的时延可能无法满足闭环控制要求；对于分钟级或小时级的运行分析、预案检索和故障复盘，MLLM 则更适合作为信息整合与解释工具。因此，可靠部署并不意味着扩大模型自治权限，而是需要根据业务时标、风险等级和校验成本划分模型可参与的环节。


\subsection{可解释性、可追溯性与责任边界}

可解释性是多模态大模型进入电力生产流程的前提。巡检识别需要说明缺陷位置和图像证据，故障诊断需要说明推理依据，调度辅助需要说明约束校验结果。若模型输出无法追溯来源，现场人员和调度人员难以判断其可信度。RAG、知识图谱和日志审计可在一定程度上提高输出可追溯性，但责任边界仍需制度化设计，尤其是当模型参与抢修方案生成或运行建议时\cite{lewis2020retrieval,deng2026repair,li2024risks}。


\subsection{评测体系与基准数据集缺失}

当前电力多模态大模型缺少统一 benchmark，不同研究在数据来源、任务定义、评价指标和场景复杂度上差异较大。单纯使用准确率或召回率难以评价模型的业务价值。面向智能电网的评测应同时覆盖识别准确性、跨场景鲁棒性、证据可追溯性、响应时延、安全违规率、幻觉率和人工采纳率。对于高风险任务，还应引入仿真环境和红队测试，评估模型在异常输入、冲突证据和规程变更情况下的表现\cite{ghafari2025comprehensive,mirshekali2025review,li2024risks}。


\subsection{工程部署与系统集成}

多模态大模型落地不仅是模型问题，也是系统工程问题。电力场景需要考虑边缘端算力、网络带宽、模型更新、业务系统接口、审计日志和安全运维流程。电力智算系统和行业大模型实践表明，算力基础设施、数据治理平台和业务流程重构是模型规模化应用的基础\cite{fangfang2025computing,du2026key}。对于巡检和应急等场景，可优先采用离线分析或人机协同方式部署；对于调度控制等高风险场景，应在仿真和沙箱环境中逐步验证。

\begin{table*}[htbp]
\centering
\heiti 表2\quad 挑战与未来方向矩阵
\vspace{-2mm}
\begin{tabular}{p{24mm}p{42mm}p{38mm}p{42mm}}
\toprule
挑战维度 & 具体表现 & 影响场景 & 未来方向 \\
\midrule
数据治理 & 数据孤岛、标注不足、质量不一 & 巡检、诊断、调度 & 标准化数据集、联邦学习、质量评估 \\
物理约束 & 生成建议可能违反潮流、限值或运行规程 & 调度控制、故障处置 & 结构化候选动作、规则校验、潮流计算、仿真验证 \\
可信推理 & 幻觉、证据冲突、来源不清 & 规程问答、应急演练 & RAG、知识图谱、证据链追踪 \\
实时部署 & 时延、算力和网络限制 & 在线监测、边缘巡检 & 模型压缩、边云协同、电力智算 \\
安全责任 & 越权操作、责任不清 & 调度、抢修、安全监管 & 人在回路、权限控制、审计机制 \\
\bottomrule
\end{tabular}
\label{tab:challenge}
\end{table*}


\section{未来研究方向}


\subsection{面向电力系统的专用多模态基础模型}

未来研究需要从通用 MLLM 调用走向电力专用多模态基础模型建设。该方向包括电力设备视觉语义库、缺陷图文对齐数据、时序--拓扑联合表示、电力规程指令集和跨场景评测集。专用模型既要提升设备与业务语义理解能力，也要避免对私有数据和高风险控制流程的不当暴露\cite{du2026key,li2024foundation_power,ouyang2026powervlm}。


\subsection{物理约束与机理模型融合}

智能电网应用不能只依赖统计相关性和语言推理。未来的关键问题不是简单地将 MLLM 与潮流计算模块相连，而是建立可审计、可复现的约束接口，使模型生成的候选方案能够被电力系统计算模块验证。图增强语言模型和配电网电压控制研究为这一方向提供了初步启示\cite{bernier2025powergraph,yang2025operator,zhangjun2023operation}，但生产级应用仍需明确模型、规则系统、仿真平台和人工调度员之间的责任边界。

第一类路径是后处理校验。MLLM 仅负责生成候选解释或候选操作，潮流计算、短路计算、暂态稳定仿真和安全规则库负责判断方案是否可执行。该路径的优点是安全边界清晰，能够最大程度保留传统电力机理模型的可验证性；其局限在于模型本身并不真正理解约束，仍可能生成大量事后被筛除的不可行方案。

第二类路径是工具调用式校验。MLLM 通过受限接口调用潮流计算、优化调度或数字孪生平台，并将工具返回的越限项、约束裕度和仿真结果纳入回答。这一路径可形成``候选建议--电力计算--结果解释''的闭环，适合调度辅助、故障处置和应急演练等需要解释性输出的场景。但其风险在于工具参数构造错误、调用权限越界以及模型对计算结果的二次解释失真，因此需要接口白名单、参数校验和调用日志审计。

第三类路径是约束感知生成。通过在训练数据、提示模板或解码约束中加入设备限值、拓扑关系、操作票规则和不可行动作清单，模型可在生成阶段减少明显违反物理约束的建议。该路径有助于降低后处理校验压力，但难以覆盖实时运行方式变化和异常工况，仍必须与外部潮流计算和仿真验证结合。总体而言，物理约束融合的目标不是让 MLLM 取代电力机理模型，而是使其在严格校验框架内承担候选方案生成、证据组织和校验结果解释等辅助功能。


\subsection{动态知识增强与持续学习}

电力规程、设备状态和运行方式会持续变化，静态知识库难以长期保持有效。未来需要动态 RAG、可更新知识图谱和持续学习机制，使模型能够处理新设备、新故障和新业务流程。同时，应建立知识版本管理和变更审计机制，确保模型引用的规程和证据处于有效状态\cite{lewis2020retrieval,ji2022survey,shen2026hgrag}。


\subsection{多智能体协同与人在回路决策}

Agent 和多智能体系统是未来趋势，而不宜作为当前智能电网 MLLM 应用的主线。ReAct 将推理和行动结合起来\cite{yao2023react}，AutoGen 等多智能体框架展示了多角色协同完成复杂任务的可能性\cite{wu2023autogen}。在电力场景中，多智能体可模拟调度员、运维人员、安监人员和知识检索模块之间的协作，用于配网故障研判和应急会商\cite{cui2026multiagent,dou2026autonomous}。但其应用必须受权限控制、审计追踪、安全隔离和人在回路确认约束。


\subsection{可信、安全与标准化评测}

多模态大模型能否进入生产级智能电网，取决于可信评测和安全标准。未来应建立覆盖缺陷识别、跨模态检索、规程问答、诊断推理和调度辅助的公开或行业级基准，评价指标应包括准确性、鲁棒性、时延、可解释性、幻觉率、安全违规率和人工采纳率。对于高风险业务，还应引入红队测试和故障注入验证，识别模型在异常输入和冲突证据下的失效模式\cite{li2024risks,ghafari2025comprehensive,mirshekali2025review}。


\section{结论}

多模态大模型为智能电网处理多源异构数据、理解电力业务语义和支撑人机协同决策提供了新的技术路径。与传统单模态模型相比，MLLM 在巡检图像理解、故障证据汇聚、规程知识问答、应急演练和调度辅助中具有更强的跨模态表达和交互能力。然而，智能电网的强安全、强实时和强物理约束属性决定了 MLLM 不能以通用生成模型形态直接进入控制闭环。未来研究应围绕电力专用多模态基础模型、知识增强可信推理、物理约束融合、标准化评测、隐私保护和人在回路机制持续推进。总体而言，多模态大模型在智能电网中的价值并非替代电力机理模型和专家体系，而是通过跨模态理解与知识增强机制，提升复杂业务场景下的感知、解释、推理和协同能力。

\vspace {5mm}
\centerline
{\zihao{5}
\heiti 参~考~文~献}

\begingroup
\renewcommand{\refname}{}
\zihao{5-}
\bibliographystyle{unsrt}
\bibliography{references}
\endgroup

\end{document}

